\subsection{正则 Noether 环的同调刻画}

定理 1 (Serre).--- 为使 Noether 局部环是正则的，必须且只需它的同调维数有限。

我们已经看到，Noether 正则局部环具有有限同调维数（命题 1）。

反之，设 A 为具有有限同调维数 n 的 Noether 局部环；由 § 3, 第 3 节, 命题 4 的推论 2 及第 5 节, 定理 1，有

\[
n = \mathrm{dh} (\mathrm{A}) = \mathrm{dp} _ {\mathrm{A}} (\kappa_ {\mathrm{A}}) = \operatorname{prof} (\mathrm{A}).
\]

若 \(n=0\)，则 A-模 \(\kappa_{A}\) 是自由的，从而 \(\mathfrak{m}_{A}=0\)，A 是域。设 \(n>0\) 并对 \(n\) 用归纳法。由于 \(\operatorname{prof}(\Lambda)>0\)，理想 \(\mathfrak{m}_{A}\) 不与 A 相伴 ( \(\S 1\) , 第 1 节, 注 2)，因而不含于 \(\mathfrak{m}_{A}^{2}\) 与 A 的相伴理想的并 (II, \(\S 1\) , 第 1 节, 命题 2)。于是 (IV, \(\S 1\) , 第 1 节, 命题 2 的推论 2)，可以找到元素 \(x\)（属于 \(\mathfrak{m}_{A}-\mathfrak{m}_{A}^{2}\)），使乘以 x 的映射 \(x_{A}\) 是单射。用 B 表示 Noether 局部环 A/ \(xA\)，并考察 A-模的序列

\[
0 \to \kappa_ {\mathrm{A}} \stackrel {i} {\longrightarrow} \mathfrak {m} _ {\Lambda} / x \mathfrak {m} _ {\Lambda} \stackrel {p} {\longrightarrow} \mathfrak {m} _ {\mathrm{B}} \to 0
\]

其中映射 \(i\) 由映射 \(a \mapsto ax\)（从 A 到 \(\mathfrak{m}_{\mathrm{A}}\)）通过取商得到，而 \(p\) 是典范满射；该序列正合。由于 \(x\) 在 \(\kappa_{\mathrm{A}}\) -向量空间 \(\mathfrak{m}_{\mathrm{A}} / \mathfrak{m}_{\mathrm{A}}^{2}\) 中的类非零，存在 A-线性映射 \(\phi : \mathfrak{m}_{\mathrm{A}} \to \kappa_{\mathrm{A}}\) 使 \(\phi(x) = 1\) ；通过取商，由 \(\phi\) 导出 \(i\) 的一个收缩，因此前面的正合列是分裂的。这蕴含关系式

\[
\mathrm{dp} _ {\mathrm{B}} \left(\mathfrak {m} _ {\mathrm{B}}\right) \leqslant \mathrm{dp} _ {\mathrm{B}} \left(\mathfrak {m} _ {\mathrm{A}} / x \mathfrak {m} _ {\mathrm{A}}\right) = \mathrm{dp} _ {\mathrm{A}} \left(\mathfrak {m} _ {\mathrm{A}}\right) <   + \infty
\]

(§ 3, 第 4 节, 命题 7 的推论 2 及 A, X, 第 135 页, 推论 1)。把同上引文的推论 2 应用于 B-模正合列 0 → mB → B → κB → 0，得到 dpB(κB) \textless{} +∞。因此环 B 具有有限同调维数 (§ 3, 第 3 节, 命题 4 的推论 2)，且深度为 n - 1 (§ 1, 第 4 节, 命题 7 及第 3 节, 命题 4 的推论)。由归纳假设，B 是正则的，从而 A 是正则的 (VIII, § 5, 第 3 节, 命题 2 的推论 1)。

因此，若 A 是 Noether 局部环，则下列三个性质等价：

\begin{enumerate}
\def\labelenumi{(\roman{enumi})}
\item
  A 是正则的；
\item
  A-模 \(\kappa_{A}\) 具有有限投射维数；
\item
  每个有限生成 A-模都具有投射维数 \(<+\infty\) 。
\end{enumerate}

定义 1.--- 称环 A 为正则的，如果它是 Noether 的，并且局部环 \(\mathrm{A}_{\mathfrak{m}}\) 都是正则的，这里 \(\mathfrak{m}\) 取遍 \(\mathrm{A}\) 的极大理想。

命题 A.--- 设 A 为 Noether 环。下列条件等价：

\begin{enumerate}
\def\labelenumi{(\roman{enumi})}
\item
  A 是正则的；
\item
  每个有限生成 A-模都具有投射维数 \(< +\infty\)；
\item
  对 A 的每个极大理想 m，A/m 的投射维数有限；
\item
  对 A 的每个素理想 p，局部环 A \(_{p}\) 都是正则的。
\end{enumerate}

设 p 为 A 的一个素理想；若 A-模 A/p 具有有限投射维数，则由 § 3, 第 2 节, 命题 3 的推论 1，\(A_{p}\) -模 \(\kappa(\mathfrak{p})\) 亦然，从而局部环 \(A_{p}\) 是正则的 (定理 1)。由此得出 (ii) 蕴含 (iv) 且 (iii) 蕴含 (i)。蕴含关系 (iv)⇒(i) 与 (ii)⇒(iii) 是显然的。

我们来证明 (i) 蕴含 (ii)。设 M 为有限生成 A-模。在假设 (i) 下，有 \(\mathrm{dp}_{\mathrm{A}_{\mathfrak{m}}}(M_{\mathfrak{m}}) \leqslant \mathrm{dh}(A_{\mathfrak{m}}) < +\infty\)（对 A 的每个极大理想 \(\mathfrak{m}\)）( \(n^{\circ} 1\) , 命题 1)；因此 M 具有有限投射维数 \(< +\infty\) ( \(\S 3, n^{\circ} 2\) , 命题 3 的推论 2)，从而得到 (ii)。

例.--- 1) 若环 A 正则，则对 A 的每个乘性子集 S，分式环 S \(^{-1}\) A 都是正则的：例如这由上面的刻画 (iii) 得出。

\begin{enumerate}
\def\labelenumi{\arabic{enumi})}
\setcounter{enumi}{1}
\item
  为使一个环正则，必须且只需它同构于一个有限正则整环族的乘积；事实上，由于正则局部环都是整环，这由每个正则环都是局部整的 (§ 1, 第 8 节) 这一事实得出。
\item
  维数 \(\leqslant 1\) 的正则整环就是 Dedekind 环 (VIII, § 5, 第 1 节, 例 1 及 VII, § 2, 第 2 节, 定理 1)。
\end{enumerate}

推论 1. 设 A 为 Noether 环。下列条件等价：

\begin{enumerate}
\def\labelenumi{(\roman{enumi})}
\item
  有 \(\mathrm{dh(A)} < +\infty\)；
\item
  A 正则且有 \(\dim(A) < +\infty\) 。
\end{enumerate}

若这些条件满足，则有 \(\dim(\Lambda) = \mathrm{d}\mathrm{h}(\mathrm{A})\) 。

若环 A 正则，则对 Λ 的每个极大理想 m 有等式 \(\dim(A_{\mathfrak{m}}) = \mathrm{dh}(A_{\mathfrak{m}})\) (第 1 节, 命题 1)，从而

\[
\mathrm{dh} (\mathrm{A}) = \sup _ {\mathfrak {m}} \mathrm{dh} (\mathrm{A} _ {\mathfrak {m}}) = \sup _ {\mathfrak {m}} \dim (\mathrm{A} _ {\mathfrak {m}}) = \dim \mathrm{A}
\]

(§ 3, 第 2 节, 命题 3 及 VIII, § 1, 第 3 节, 命题 8)。另一方面，若 dh(Λ) \textless{} +∞，则由命题 4，环 A 正则。推论得证。

存在维数无限的正则 Noether 环 (VIII, § 5, 习题 6)。

推论 2. 正则环是正规的、Gorenstein 的和 Macaulay 的。

事实上，正则局部环是整闭的 (VIII, § 5, 第 2 节, 定理 1 的推论 1)、Gorenstein 的 (§ 3, 第 9 节, 例 4) 和 Macaulay 的 (§ 2, 第 5 节, 例 6)。

推论 3.--- 设 A 为 Noether 环，J 为 A 的理想，Â 为 A 关于 J-进拓扑的分离完备化。

\begin{enumerate}
\def\labelenumi{\alph{enumi})}
\item
  为使环 \(\widehat{\mathbf{A}}\) 正则，必须且只需对每个包含 J 的极大理想 \(\mathfrak{m}\)（属于 \(\Lambda\)），环 \(\mathrm{A}_{\mathfrak{m}}\) 都是正则的。
\item
  若环 A 正则，则环 \(\hat{\Lambda}\) 正则。若环 \(\hat{\mathrm{A}}\) 正则且理想 J 含于 A 的根中，则环 \(\Lambda\) 正则。
\end{enumerate}

由 III, § 3, 第 4 节的命题 8，为使环 \(\widehat{\Lambda}\) 正则，必须且只需 \(\widehat{A}_{\widehat{\mathfrak{m}}}\) 都正则（对 A 的每个包含 J 的极大理想 \(\mathfrak{m}\)）。由于局部环 \(\widehat{A}_{\widehat{\mathfrak{m}}}\) 与 \(A_{\mathfrak{m}}\) 的完备化同构（同上引文），断言 a) 由 VIII, § 5, 第 1 节, 命题 1 的推论得出。断言 b) 由 a) 得出。

推论 4.--- 设 A 为正则环，P 为有限生成投射 A-模。对称代数 \(\mathsf{S}_{\mathrm{A}}(\mathrm{P})\) 是正则环。

设 p 为 \(\mathbf{S}_{\mathrm{A}}(\mathrm{P})\) 的一个素理想，q 为它在 A 中的原像。局部环 \(\mathbf{S}_{\mathrm{A}}(\mathrm{P})_{\mathfrak{p}}\) 是环 \(\mathbf{S}_{\Lambda}(\mathrm{P})_{\mathfrak{q}}\) 的一个分式环，而后者同构于 \(\mathbf{S}_{\mathrm{A}_{q}}(\mathrm{P}_{\mathfrak{q}})\) (A, III, 第 72 页, 命题 7)；只需证明后者是正则的。于是归结为 A 是局部环的情形；而这时 P 是有限型的自由模。由第 1 节的命题 1 及 A, X, 第 143 页, 推论 1，有 \(\mathrm{dh}(\mathbf{S}_{\mathrm{A}}(\mathrm{P})) = \mathrm{dh}(\mathrm{A}) + \mathrm{rg}_{\mathrm{A}}(\mathrm{P}) < +\infty\)，从而 \(\mathbf{S}_{\mathrm{A}}(\mathrm{P})\) 由命题 4 是正则的。

推论 5. 设 A 为正则环，\((\mathrm{T}_i)_{i\in I}\) 为有限的未定元族。多项式环 \(\mathrm{A}[(\mathrm{T}_i)_{i\in I}]\) 与形式幂级数环 \(\mathrm{A}[(\mathrm{T}_i)_{i\in I}]\) 都是正则的。

这由推论 4 及推论 3 b) 得出。

